\chapter{Tularemia}
\index{Tularemia}

\section*{Definition and Overview}
Tularemia is a highly infectious, potentially life-threatening zoonotic disease caused by the bacterium \textit{Francisella tularensis}, a small, Gram-negative, facultative intracellular coccobacillus. Commonly referred to by several historical, geographic, or occupational terms such as ``rabbit fever,'' ``deer fly fever,'' ``Francis' disease,'' or ``Ohara's disease,'' this pathogen is characterized by its extraordinary virulence. A remarkably low infectious dose is required to establish infection; inhalation, ingestion, or inoculation through minor skin abrasions of as few as 10 to 50 bacterial cells can lead to severe systemic illness in humans. Because of this high infectivity, its potential for aerosol dissemination, and the severity of the clinical disease it causes, \textit{Francisella tularensis} is classified globally as a Tier 1 Select Agent and is recognized by public health and national defense organizations as a potential agent of biological warfare and bioterrorism.

The clinical presentation of tularemia is highly diverse, with the disease manifesting in several distinct forms based primarily on the route of exposure. These include ulceroglandular, glandular, oculoglandular, oropharyngeal, pneumonic, and typhoidal forms. Historically first described in the early 20th century in Tulare County, California, by Dr. Edward Francis, the disease remains a significant public health focus due to its persistence in wild animal reservoirs, its complex transmission cycles involving a broad range of mammalian hosts and arthropod vectors, and the critical need for rapid diagnosis and appropriate antibiotic intervention. Understanding the pathophysiology, clinical indicators, and therapeutic pathways of tularemia is essential for clinicians, public health officials, and those who work or recreate in endemic outdoor environments.

\section*{Epidemiology}
The epidemiology of tularemia is complex, reflecting the pathogen's ability to infect more than 250 species of wild and domestic animals, including mammals, birds, reptiles, fish, and various invertebrates. The disease is endemic throughout the Northern Hemisphere, with cases reported across North America, continental Europe, northern Asia, and parts of the Middle East. It is absent in the Southern Hemisphere, except for rare, isolated reports.

In North America, the disease is primarily caused by two subspecies of the bacterium: \textit{Francisella tularensis} subsp. \textit{tularensis} (Type A) and \textit{Francisella tularensis} subsp. \textit{holarctica} (Type B). Type A is highly virulent in humans, frequently associated with land rabbits (such as cottontails and jackrabbits) and tick vectors, and is primarily restricted to North America. Type B is generally less virulent, associated with aquatic environments, water rats, muskrats, and beavers, and is widespread across North America, Europe, and Asia.

In the United States, tularemia is a relatively rare but persistent public health issue, with approximately 100 to 300 cases reported annually. The disease is highly concentrated in specific geographic regions, particularly in the south-central and western states, including Arkansas, Oklahoma, Missouri, Kansas, and South Dakota. In Europe, countries like Sweden, Finland, Norway, and Turkey report periodic outbreaks, which are frequently linked to mosquito bites, agricultural activities, or contact with contaminated natural water sources.

Demographically, tularemia exhibits a strong occupational and recreational bias. Historically, individuals who engage in activities that bring them into contact with wildlife or vectors are at the highest risk. This includes hunters, trappers, laboratory personnel, veterinarians, agricultural workers, and landscapers. Recreationally, hikers, campers, and home gardeners are also at risk. The disease displays a clear seasonal distribution, with peak incidence in the late spring and summer months due to increased tick and deer fly activity, and a secondary peak in winter, corresponding to hunting seasons for rabbits and other small game. Males are disproportionately affected, representing approximately 70\% of all reported cases, a disparity attributed to higher occupational and recreational exposure.

\section*{Aetiology and Pathophysiology}
The underlying cause of tularemia is infection with the bacterium \textit{Francisella tularensis}, a small, non-motile, pleomorphic, Gram-negative coccobacillus. The organism is a facultative intracellular pathogen, meaning it can survive and replicate both inside and outside host cells, with a marked preference for replicating within host macrophages.
\begin{itemize}
    \item \textbf{Causative Organism}: \textit{Francisella tularensis} is a fastidious bacterium that requires cysteine-supplemented media (such as cystine heart agar or chocolate agar) for laboratory isolation. The bacterium is surrounded by a lipid-rich capsule that plays a critical role in its virulence by preventing phagocytosis, inhibiting complement-mediated lysis, and suppressing early inflammatory responses in the host.
    \item \textbf{Subspecies Differentiation}: There are four recognized subspecies of \textit{Francisella tularensis}, but only two cause significant human disease. Subspecies \textit{tularensis} (Type A) is highly virulent, with an untreated mortality rate of up to 30\% in pulmonary cases. Subspecies \textit{holarctica} (Type B) is less virulent, rarely causing death even if untreated, and is often linked to waterborne transmission.
    \item \textbf{Transmission Pathways}: Transmission to humans occurs through several distinct routes:
    \begin{itemize}
        \item Vector bites: Transmission occurs via the bites of infected hard ticks (e.g., \textit{Dermacentor variabilis}, \textit{Dermacentor andersoni}, and \textit{Amblyomma americanum}) or biting insects, such as deer flies (\textit{Chrysops} species) and, in Europe, mosquitoes (\textit{Aedes} and \textit{Anopheles} species).
        \item Direct animal contact: Inoculation through minor skin abrasions or mucous membranes when handling infected animal tissues, fluids, or pelts. This is common among hunters skinning rabbits or muskrats.
        \item Ingestion: Consuming undercooked meat of infected animals or drinking water contaminated by the excretions or carcasses of infected rodents.
        \item Inhalation: Inhaling aerosolized bacteria during agricultural tasks (e.g., mowing lawns, clearing brush, or threshing grain where animal carcasses or vector feces are present) or in laboratory accidents.
        \item Domestic animal transmission: Scratches or bites from domestic cats that have hunted and killed infected rabbits or rodents.
    \end{itemize}
    \item \textbf{Intracellular Replication Loop}: Once inside the host, the bacteria are engulfed by macrophages. Unlike most bacteria, which are destroyed within the phagolysosome, \textit{F. tularensis} escapes the phagosome into the host cell's cytoplasm. It does this by suppressing the assembly of the NADPH oxidase complex, preventing the respiratory burst. In the cytoplasm, the bacteria replicate rapidly, eventually causing host cell lysis and releasing progeny to infect neighboring cells.
    \item \textbf{Lymphatic and Systemic Dissemination}: The bacteria travel via the lymphatic system from the entry site to regional lymph nodes, leading to severe inflammation, focal necrosis, and painful lymphadenitis. If the infection is not contained by regional lymph nodes, the bacteria enter the bloodstream (bacteraemia), allowing dissemination to systemic organs such as the liver, spleen, lungs, and bone marrow. This systemic spread leads to the formation of granulomatous lesions with central caseating necrosis, pathologically resembling tuberculosis.
\end{itemize}

\section*{Clinical Features}
The clinical presentation of tularemia is highly variable and depends on the route of bacterial entry into the body. The incubation period typically ranges from 3 to 5 days, but can extend from 1 to 14 days. The disease begins abruptly with non-specific, systemic symptoms including high fever (often reaching 40 degrees Celsius), chills, generalized body aches (myalgia), severe headache, joint pain (arthralgia), sore throat, and profound fatigue. Following this initial prodrome, the disease resolves into one of six classical clinical forms:
\begin{itemize}
    \item \textbf{Ulceroglandular Tularemia}: This is the most common form, representing 75\% to 85\% of all cases. It occurs after skin exposure via arthropod bites or direct contact with infected tissue. It is characterized by:
    \begin{itemize}
        \item A single, painful skin lesion at the site of inoculation. The lesion begins as a red papule, progresses to a pustule, and then ulcerates. The resulting ulcer is typically circular, punched-out, has a necrotic base, and is surrounded by a raised, hardened (indurated) border.
        \item Marked, painful regional lymphadenopathy. The lymph nodes draining the site of the ulcer (typically inguinal, axillary, or cervical) swell significantly and are exquisitely tender. These swollen nodes, or buboes, may fluctuate, become fluctuant, and occasionally suppurate (drain pus) if untreated.
    \end{itemize}
    \item \textbf{Glandular Tularemia}: This form occurs through the same routes as the ulceroglandular form but lacks a visible skin ulcer at the portal of entry. Patients present with painful, swollen regional lymph nodes and systemic febrile symptoms.
    \item \textbf{Oculoglandular Tularemia}: This form occurs when the bacterium is introduced directly into the eye, often via contaminated fingers, splashes of body fluids from infected animals, or dust. It accounts for 1\% to 5\% of cases and is characterized by:
    \begin{itemize}
        \item Severe, unilateral conjunctivitis with marked eye pain, irritation, and foreign body sensation.
        \item Multiple small, yellowish conjunctival nodules or shallow ulcers.
        \item Exquisite tenderness and swelling of the preauricular, submandibular, and cervical lymph nodes on the affected side.
        \item Photophobia (light sensitivity) and watery or purulent discharge from the eye.
    \end{itemize}
    \item \textbf{Oropharyngeal Tularemia}: Acquired by ingestion of contaminated water, undercooked wild game (especially rabbit), or through aerosol deposition in the mouth. It is characterized by:
    \begin{itemize}
        \item Severe pharyngitis or tonsillitis, often presenting with a grayish-white exudative pseudomembrane that can mimic diphtheria.
        \item Significant bilateral cervical, submandibular, or retropharyngeal lymphadenopathy.
        \item Abdominal pain, nausea, vomiting, and diarrhoea, which can lead to rapid dehydration.
    \end{itemize}
    \item \textbf{Pneumonic Tularemia}: This is the most severe, life-threatening form of the disease. It can occur as a primary infection following the inhalation of aerosolized bacteria (e.g., during agricultural activities or laboratory accidents) or secondarily from systemic haematogenous dissemination. Symptoms include:
    \begin{itemize}
        \item Acute onset of high fever, chills, and progressive dyspnoea (shortness of breath).
        \item Dry, non-productive cough, which may become productive of purulent or blood-tinged sputum (haemoptysis).
        \item Pleuritic chest pain (pain that worsens with breathing) and chest tightness.
        \item Physical findings of consolidation, and radiographic evidence of atypical pneumonia, lobar infiltrates, pleural effusion, or hilar lymphadenopathy.
    \end{itemize}
    \item \textbf{Typhoidal Tularemia}: This is a severe, systemic form of the disease that presents without localized signs of infection (such as skin ulcers or localized lymphadenopathy). It occurs in 5\% to 15\% of cases and is characterized by:
    \begin{itemize}
        \item Persistent high fever, profound weakness, weight loss, and altered mental status.
        \item Abdominal pain, hepatosplenomegaly (enlarged liver and spleen), and diarrhoea.
        \item High risk of progression to septic shock, acute respiratory distress syndrome (ARDS), and multi-organ failure. It carries the highest mortality rate of all forms of tularemia due to delayed diagnosis.
    \end{itemize}
\end{itemize}

\section*{Differential Diagnosis}
Because the clinical features of tularemia are diverse and often mimic other diseases, clinicians must consider a wide range of alternative diagnoses based on the presenting form.
\begin{itemize}
    \item \textbf{Plague (\textit{Yersinia pestis})}: Plague is a primary consideration in patients presenting with acute, painful lymphadenopathy (buboes). Plague is typically more rapid in onset and progression, is associated with flea bites or exposure to wild rodents in specific endemic regions, and is characterized by a higher degree of systemic toxicity and a higher incidence of purpuric skin lesions (``black death'').
    \item \textbf{Cat Scratch Disease (\textit{Bartonella henselae})}: This bacterial infection is a common cause of subacute regional lymphadenopathy, especially in children and young adults with a history of kitten scratches or bites. While it can cause skin papules and swollen nodes (typically axillary or cervical), it is generally a milder, self-limiting disease with less severe systemic symptoms (fever is usually mild or absent).
    \item \textbf{Tuberculosis (\textit{Mycobacterium tuberculosis})}: Tuberculosis can present with chronic lymphadenitis (scrofula) or pulmonary symptoms that mimic pneumonic tularemia. However, tuberculosis typically has a much slower, more indolent progression, characteristic chest radiographic findings (e.g., upper lobe cavitary lesions), and is diagnosed via sputum smear for acid-fast bacilli, mycobacterial culture, or interferon-gamma release assays.
    \item \textbf{Cutaneous Anthrax (\textit{Bacillus anthracis})}: This must be differentiated from ulceroglandular tularemia. Cutaneous anthrax is characterized by a painless, pruritic papule that quickly develops into a central black eschar surrounded by marked, gelatinous non-pitting oedema. In contrast, the skin ulcer in tularemia is highly painful and lacks the classic black eschar.
    \item \textbf{Sporotrichosis (\textit{Sporothrix schenckii})}: This fungal infection, often associated with minor skin trauma from rose thorns or soil, causes cutaneous nodules that spread along lymphatic pathways (sporotrichoid spread). Unlike tularemia, sporotrichosis is typically localized, chronic, and lacks acute systemic features like high fever and severe constitutional symptoms.
    \item \textbf{Atypical Pneumonia (e.g., \textit{Mycoplasma pneumoniae}, \textit{Legionella pneumophila}, Q Fever)}: These pathogens cause acute pulmonary infections that are clinically and radiographically indistinguishable from pneumonic tularemia. Differential diagnosis depends on detailed exposure histories (such as exposure to livestock for Q fever, water systems for Legionella) and specific serological or molecular tests.
    \item \textbf{Brucellosis (\textit{Brucella} species)}: Brucellosis causes systemic symptoms such as undulating fever, sweats, joint pain, and fatigue, which overlap with typhoidal tularemia. Brucellosis is distinguished by its more subacute or chronic course, association with consuming unpasteurized dairy products or occupational exposure to livestock, and specific serology.
\end{itemize}

\section*{When to Seek Medical Care}
Tularemia is a serious disease that can deteriorate rapidly without appropriate treatment. It is critical to seek medical attention if you experience a sudden onset of high fever, chills, unexplained swollen lymph nodes, or a painful, non-healing skin sore. This is particularly important if you have a history of recent outdoor activities, tick or insect bites, contact with wild animals (especially rabbits and rodents), or if you live in or have recently visited an area where tularemia is known to occur.

For life-threatening symptoms, immediate emergency medical intervention is required. Call your local emergency services (e.g., \textbf{911} in the United States, \textbf{000} in many countries, \textbf{112} in Europe, or \textbf{999} in the United Kingdom) or go to the nearest emergency department if you experience any of the following:
\begin{itemize}
    \item Severe difficulty breathing, short breath, or chest pain.
    \item High, persistent fever accompanied by confusion, severe lethargy, or inability to wake up.
    \item Signs of shock, such as cold, clammy, or pale skin, rapid heart rate, shallow breathing, and dizziness or fainting.
    \item Bluish discoloration of the lips, face, or fingernails (cyanosis).
    \item Sudden, severe swelling of the throat that interferes with swallowing or breathing.
\end{itemize}

\section*{Diagnosis and Investigations}
Diagnosing tularemia requires a high degree of clinical suspicion, particularly regarding exposure history, as the disease is rare and its clinical manifestations are non-specific.
\begin{itemize}
    \item \textbf{Clinical History and Physical Examination}: The diagnostic process begins with a detailed history, including recent travel, outdoor activities, tick or insect bites, contact with wild or domestic animals, and occupation. The physical examination focuses on identifying skin ulcers, evaluating the distribution and tenderness of swollen lymph nodes, and assessing respiratory and ocular findings.
    \item \textbf{Serology}: The most widely used diagnostic method is serological testing to detect antibodies against \textit{Francisella tularensis}.
    \begin{itemize}
        \item Tube Agglutination and Microagglutination: These are the standard serological tests. A single titre of 1:160 or higher is suggestive of infection, but a definitive diagnosis requires a fourfold increase in antibody titre between acute-phase serum (collected within the first week of illness) and convalescent-phase serum (collected 2 to 4 weeks later).
        \item Enzyme-Linked Immunosorbent Assay (ELISA): ELISA tests can detect IgM and IgG antibodies, offering a rapid and sensitive diagnostic option.
    \end{itemize}
    \item \textbf{Polymerase Chain Reaction (PCR)}: PCR assays can detect the DNA of \textit{Francisella tularensis} directly from clinical specimens such as ulcer swabs, lymph node aspirates, sputum, or whole blood. PCR is highly sensitive and specific, providing rapid results early in the course of the disease before antibodies have developed.
    \item \textbf{Microbiological Culture and Isolation}: Culturing \textit{F. tularensis} from clinical specimens (blood, lymph node aspirates, pleural fluid, or ulcer exudate) provides a definitive diagnosis. However, \textit{F. tularensis} is extremely infectious, and laboratory personnel are at high risk of acquiring the infection through aerosols.
    \item \textbf{Biosafety Warnings}: Because of the extreme risk of laboratory transmission, clinicians must notify laboratory staff immediately if tularemia is suspected. Specimens must be handled under Biosafety Level 3 (BSL-3) containment conditions. Routine laboratory processing should be avoided.
    \item \textbf{Imaging and Other Tests}: For patients with suspected pulmonary involvement, a chest X-ray or CT scan should be performed. Typical findings include patchy or lobar infiltrates, bronchopneumonia, pleural effusion, or hilar lymphadenopathy. Basic blood tests may show non-specific signs of infection, such as an elevated white blood cell count (leukocytosis), elevated C-reactive protein (CRP), and elevated erythrocyte sedimentation rate (ESR).
\end{itemize}

\section*{Management}
The primary treatment for tularemia is timely and appropriate antibiotic therapy. Because the disease can progress rapidly and cause severe complications, treatment should be initiated empirically if clinical suspicion is high, without waiting for laboratory confirmation.
\begin{itemize}
    \item \textbf{First-Line Antibiotics (Aminoglycosides)}: For severe or systemic cases of tularemia (such as pneumonic or typhoidal forms), aminoglycosides are the preferred treatment:
    \begin{itemize}
        \item Streptomycin: Historically considered the drug of choice, it is highly bactericidal. It is administered intramuscularly for 7 to 10 days. However, its availability is limited in many countries.
        \item Gentamicin: Widely available and equally effective, gentamicin is the standard intravenous alternative. It is administered for 7 to 10 days, with blood levels monitored to prevent toxicity.
    \end{itemize}
    \item \textbf{Alternative Antibiotics (Tetracyclines and Fluoroquinolones)}: For mild to moderate cases (such as localized ulceroglandular or glandular forms), oral antibiotics are appropriate:
    \begin{itemize}
        \item Doxycycline: An effective oral treatment, but because it is bacteriostatic, it must be taken for a longer duration (14 to 21 days) to prevent relapse.
        \item Ciprofloxacin: A fluoroquinolone that is highly effective due to its excellent intracellular penetration and bactericidal activity. The recommended course is 10 to 14 days.
    \end{itemize}
    \item \textbf{Surgical Intervention}: Swollen lymph nodes (buboes) that become fluctuant and painful despite antibiotic therapy may require needle aspiration or surgical drainage to relieve pain and promote healing. Excision is generally avoided as it can spread the infection.
    \item \textbf{Supportive Care}: Patients with severe disease, respiratory distress, or septic shock require intensive supportive care in a hospital setting:
    \begin{itemize}
        \item Intravenous fluids to maintain hemodynamic stability.
        \item Supplemental oxygen or mechanical ventilation for respiratory failure.
        \item Vasopressors to support blood pressure in septic shock.
        \item Adequate pain management and wound care for skin ulcers.
    \end{itemize}
    \item \textbf{Post-Exposure Prophylaxis (PEP)}: In cases of confirmed or highly probable exposure (such as a laboratory accident or a bioterrorism event), post-exposure prophylaxis with oral doxycycline or ciprofloxacin for 14 days is recommended to prevent the onset of disease.
\end{itemize}

\section*{Complications}
If left untreated, or if treatment is delayed, tularemia can cause severe, life-threatening complications due to the spread of the bacteria through the bloodstream.
\begin{itemize}
    \item \textbf{Acute Respiratory Distress Syndrome (ARDS)}: A severe complication of pneumonic tularemia where fluid accumulates in the lungs' air sacs, leading to severe oxygen deprivation and respiratory failure requiring mechanical ventilation.
    \item \textbf{Septic Shock and Multi-Organ Failure}: Systemic spread of the infection (septicemia) can trigger a massive inflammatory response, leading to a dangerous drop in blood pressure, poor tissue perfusion, and failure of critical organs such as the kidneys, liver, and heart.
    \item \textbf{Meningitis}: In rare instances, the bacteria can spread to the central nervous system, causing inflammation of the membranes surrounding the brain and spinal cord. Tularemic meningitis is associated with high mortality and long-term neurological deficits if not treated promptly.
    \item \textbf{Pericarditis and Endocarditis}: The infection can spread to the heart, causing inflammation of the surrounding sac (pericarditis) or the inner lining of the heart chambers and valves (endocarditis). This can lead to chest pain, fluid accumulation around the heart (pericardial effusion), or heart valve damage.
    \item \textbf{Osteomyelitis and Septic Arthritis}: Hematogenous spread can lead to infections in the bones (osteomyelitis) or joints (septic arthritis), causing severe pain, swelling, and potential joint destruction.
    \item \textbf{Renal Failure}: Sepsis-induced kidney injury or direct bacterial damage can cause acute kidney injury, which may require temporary dialysis.
    \item \textbf{Treatment-Related Complications}: The antibiotics used to treat tularemia have potential side effects. Aminoglycosides (streptomycin and gentamicin) are associated with nephrotoxicity (kidney damage) and ototoxicity (hearing loss and balance problems). Doxycycline can cause severe gastrointestinal irritation and photosensitivity, while ciprofloxacin carries risks of tendonitis, tendon rupture, and central nervous system side effects.
\end{itemize}

\section*{Prognosis and Outcomes}
The prognosis for individuals with tularemia is highly dependent on the subspecies of the infecting organism, the clinical form of the disease, the timing of antibiotic initiation, and the patient's underlying health status.
With early and appropriate antibiotic therapy, the prognosis is generally excellent. The overall mortality rate for treated cases of tularemia is less than 2\%. Most patients recover completely, though the recovery process can be slow. Ulcers typically heal over several weeks, and swollen lymph nodes may take weeks to months to return to their normal size.

If left untreated, the prognosis is much worse, particularly for infections caused by the highly virulent Type A subspecies. The mortality rate for untreated Type A ulceroglandular tularemia ranges from 5\% to 15\%, while the mortality rate for untreated pneumonic or typhoidal tularemia can exceed 30\% to 60\%. In contrast, untreated Type B infections (common in Europe and Asia) are rarely fatal, though they can cause prolonged, debilitating illness.

Factors that negatively impact prognosis include:
\begin{itemize}
    \item Delay in seeking medical care or starting appropriate antibiotic therapy.
    \item Infection with \textit{Francisella tularensis} subsp. \textit{tularensis} (Type A).
    \item Development of severe complications such as pneumonic tularemia, ARDS, or septic shock.
    \item Advanced age or pre-existing immunocompromised states (e.g., HIV/AIDS, active cancer treatment, or immunosuppressive therapy).
\end{itemize}
Long-term immunity typically follows recovery from tularemia, though rare cases of reinfection have been reported.

\section*{Prevention and Self-Care}
Preventing tularemia relies primarily on avoiding exposure to the bacteria, protecting against vector bites, and practicing safe handling of animals. There is currently no widely available vaccine for the general public, although an investigational live vaccine has been used to protect laboratory personnel.
\begin{itemize}
    \item \textbf{Protection Against Vectors}:
    \begin{itemize}
        \item Use insect repellents containing DEET, picaridin, or IR3535 on exposed skin and clothing when spending time outdoors in endemic areas.
        \item Wear long-sleeved shirts, long trousers, and tuck pants into socks to minimize skin exposure to ticks and flies.
        \item Perform frequent tick checks on yourself, children, and pets after outdoor activities, and remove ticks promptly using fine-tipped tweezers by pulling straight out.
    \end{itemize}
    \item \textbf{Safe Animal Handling}:
    \begin{itemize}
        \item Avoid handling dead or sick wild animals. If you must handle wildlife, wear thick, impervious gloves (such as rubber or vinyl gloves) and protective eyewear.
        \item Wash hands thoroughly with soap and warm water after handling any animal or carcass.
        \item Cook wild game meat (especially rabbit and squirrel) thoroughly to an internal temperature of at least 74 degrees Celsius (165 degrees Fahrenheit) to destroy any bacteria.
    \end{itemize}
    \item \textbf{Landscaping and Agricultural Precautions}:
    \begin{itemize}
        \item Check lawns and fields for animal carcasses before mowing or brush-clearing.
        \item Wear a high-quality protective mask (such as an N95 respirator) when mowing lawns, clearing brush, or handling hay in endemic areas to prevent inhaling aerosolized bacteria from carcasses or contaminated soil.
    \end{itemize}
    \item \textbf{Water Safety}:
    \begin{itemize}
        \item Avoid drinking untreated water from lakes, rivers, or shallow wells. Ensure all drinking water is properly treated, filtered, or boiled, especially in areas where Type B tularemia is endemic.
    \end{itemize}
    \item \textbf{Pet Care}:
    \begin{itemize}
        \item Keep domestic pets, particularly cats and dogs, indoors or under close supervision to prevent them from hunting wild rodents or rabbits.
        \item Apply regular tick and flea prevention treatments to pets.
        \item Consult a veterinarian immediately if a pet shows signs of illness (such as fever, lethargy, or swollen lymph nodes) after contact with wildlife.
    \end{itemize}
\end{itemize}

\section*{Sources and Further Reading}
\begin{itemize}
    \item \textbf{Centers for Disease Control and Prevention (CDC)}: Offers comprehensive guidelines on the symptoms, transmission, diagnosis, treatment, and prevention of tularemia. URL: \url{https://www.cdc.gov/tularemia/index.html}
    \item \textbf{World Health Organization (WHO)}: Provides international epidemiological data, disease outbreak news, and global biosafety manuals covering \textit{Francisella tularensis}. URL: \url{https://www.who.int/health-topics/tularemia}
    \item \textbf{European Centre for Disease Prevention and Control (ECDC)}: Features factsheets, surveillance reports, and distribution maps of tularemia cases across European countries. URL: \url{https://www.ecdc.europa.eu/en/tularemia}
    \item \textbf{National Institutes of Health (NIH) - National Institute of Allergy and Infectious Diseases (NIAID)}: Details ongoing scientific research, clinical trials, and vaccine development efforts regarding tularemia and biodefense pathogens. URL: \url{https://www.niaid.nih.gov/diseases-conditions/tularemia}
    \item \textbf{Public Health Agency of Canada (PHAC)}: Publishes pathogen safety data sheets (PSDS) and infectious disease guidelines for laboratory workers and the public regarding \textit{Francisella tularensis}. URL: \url{https://www.canada.ca/en/public-health/services/diseases/tularemia.html}
    \item \textbf{The Merck Manual Professional Edition}: Provides a highly detailed medical reference guide on the etiology, pathophysiology, symptoms, diagnosis, and treatment protocols for all clinical forms of tularemia. URL: \url{https://www.merckmanuals.com/professional/infectious-diseases/gram-negative-bacilli/tularemia}
\end{itemize}